\section{The normally convergent Gauss--Hurwitz difference}\label{sec:subtraction}

\subsection{The summand and its exact summation}
Define
\begin{equation}\label{eq:rn}
 r_n(u)=\frac{\Gamma(n+a)\Gamma(n+b)}
 {\Gamma(n+1)\Gamma(n+a+b+u)},\qquad
 G(u)=\frac{\Gamma(a)\Gamma(b)\Gamma(u)}
 {\Gamma(a+u)\Gamma(b+u)}.
\end{equation}
For $\Rea u>0$ one has
\begin{equation}\label{eq:gauss}
 \sum_{n=0}^\infty r_n(u)=G(u).
\end{equation}
Here is a direct proof, which also fixes the normalization.  The beta
integral gives
\[
 r_n(u)=\frac{\Gamma(n+a)}{\Gamma(n+1)\Gamma(a+u)}
 \int_0^1 t^{n+b-1}(1-t)^{a+u-1}\,\dd t.
\]
Summing $\Gamma(n+a)t^n/n!=\Gamma(a)(1-t)^{-a}$ gives
$\Gamma(a)B(b,u)/\Gamma(a+u)=G(u)$.
For real positive parameters Tonelli's theorem applies; on compact subsets
of $\Rea u>0$ the same calculation follows by absolute convergence and
analytic continuation.  This is Gauss's classical summation
\cite[15.4.20]{dlmf-gauss}.

\subsection{Explicit coefficient polynomials}
Let $B_k(y)$ denote the Bernoulli polynomial, with
$B_1(y)=y-\tfrac12$.  For $k\ge1$ put
\begin{equation}\label{eq:ell}
 \ell_k(u;c)=\frac{(-1)^{k+1}}{k(k+1)}
 \left\{
 B_{k+1}(a-c)+B_{k+1}(b-c)-B_{k+1}(1-c)
 -B_{k+1}(a+b+u-c)\right\}.
\end{equation}
Define $A_0=1$ and, successively,
\begin{equation}\label{eq:Arec}
 j A_j(u;c)=\sum_{k=1}^j k\ell_k(u;c)A_{j-k}(u;c).
\end{equation}
These are polynomials with rational coefficients in $a,b,u,c$.
The first nontrivial one is
\begin{equation}\label{eq:A1}
 A_1(u;c)=c-ab+(c-a-b+\tfrac12)u-\tfrac12u^2.
\end{equation}
The recurrence is a finite exact formula at every prescribed order; no
infinite asymptotic expansion is numerically summed.

\begin{lemma}[Uniform finite-order expansion]\label{lem:asymptotic}
For any fixed $K\ge1$,
\begin{equation}\label{eq:asymptotic}
 r_n(u)=(n+c)^{-1-u}
 \left\{\sum_{j=0}^{K-1}\frac{A_j(u;c)}{(n+c)^j}
 +O\bigl((n+c)^{-K}\bigr)\right\}.
\end{equation}
The error is uniform on compact subsets of the stated parameter domain,
including compact sets of complex $u$.  Arbitrary fixed-order parameter
and spectral derivatives have the corresponding uniform expansions;
spectral differentiation introduces at most fixed powers of $\log(n+c)$.
\end{lemma}
\begin{proof}
Apply the finite Stirling expansion for $\log\Gamma(X+d)$, with
$X=n+c$, to the two numerator and two denominator factors.  The common
$X\log X-X$ and constant terms cancel.  The coefficient of $\log X$ is
$-1-u$, and the inverse-power coefficients are exactly \eqref{eq:ell}.
Exponentiating a finite expansion with its remainder gives
\eqref{eq:asymptotic}.  Differentiating the formal exponential yields
\eqref{eq:Arec}.  Uniformity follows from the uniform remainder on compact
sets of shifts in the Stirling expansion \cite[\S5.11]{dlmf-gamma}.
For a fully analytic justification of derivatives, enlarge a given compact
parameter set slightly, apply the uniform expansion there, and use Cauchy's
integral formula for the remainder.  Only the explicit factor
$X^{-u}$ contributes the possible logarithmic powers.  Finitely many small
$n$ are handled separately by entire reciprocal gamma.
\end{proof}

\begin{theorem}[Gauss--Hurwitz subtraction]\label{thm:subtraction}
For $K\ge1$, the series
\begin{equation}\label{eq:RK}
 R_K(u;c)=\sum_{n=0}^\infty
 \left\{r_n(u)-\sum_{j=0}^{K-1}
       \frac{A_j(u;c)}{(n+c)^{1+u+j}}\right\}
\end{equation}
converges absolutely and locally normally on $\Rea u>-K$.
It is holomorphic there, and
\begin{equation}\label{eq:master}
 \boxed{R_K(u;c)=G(u)-\sum_{j=0}^{K-1}
 A_j(u;c)\zeta(1+u+j,c).}
\end{equation}
Every apparent pole of the right side inside this half-plane is removable.
Arbitrary fixed-order spectral and parameter derivatives can be taken
termwise in the \emph{bracketed} series.
\end{theorem}
\begin{proof}
On a compact subset with $\Rea u\ge-K+\delta$, Lemma~\ref{lem:asymptotic}
bounds the bracket by a constant times $(n+c)^{-1-\delta}$; derivatives
add only a fixed logarithmic factor.  These majorants are summable.
The Weierstrass theorem and the differentiated version of the same estimate
prove normal convergence and holomorphy.  In $\Rea u>0$, both separate
series converge and \eqref{eq:master} follows from \eqref{eq:gauss} and
the defining Dirichlet series for Hurwitz zeta.  The identity theorem gives
\eqref{eq:master} on the larger domain, with all singularities removable
because its left side is holomorphic.
\end{proof}

\begin{remark}[Ordinary sums, not unassigned divergent sums]
At $u=-N$, $N<K$, the two components inside the braces in
\eqref{eq:RK} usually have divergent sums separately.  The theorem evaluates
the single absolutely convergent sum of their difference.  It does not
assign values to either divergent series or authorize rearranging the
separate pieces.  The value depends on the stated reference shift $c$;
Section~\ref{sec:primitives} gives its exact dependence.
\end{remark}

Increasing the number of counterterms gives the exact relation
\begin{equation}\label{eq:Krec}
 R_{K+1}(u;c)=R_K(u;c)-A_K(u;c)\zeta(1+u+K,c),
 \qquad \Rea u>-K.
\end{equation}
Thus additional subtractions improve summability without creating a new
unknown constant.  The meromorphic continuation of this relation remains
valid wherever its combined terms are defined.
